
\thispagestyle{cover}%
\newgeometry{bottom=25mm,left=20mm,top=15mm,right=20mm}
\hspace{-47pt}
\begin{minipage}[l]{0.2\columnwidth}
\vspace{6mm}
\includegraphics[width=0.8\columnwidth]{tpl/Sesame_logo.png}\\
\end{minipage}
\hfill
\begin{minipage}[l]{0.5\columnwidth}
\centering
\footnotesize
\textbf{{République Tunisienne}}\\
\vspace{1.5mm}
\textbf{{Ministère de l'Enseignement Supérieur\\
et de la Recherche Scientifique}}\\
\vspace{1.5mm}
%\textbf{{Université de Tunis El Manar}}\\
%\vspace{1.5mm}
%\textbf{{Institut Supérieur d'Informatique d’El Manar}}
\end{minipage}
\hfill
\begin{minipage}[l]{0.02\columnwidth}
\end{minipage}
\hfill
\begin{minipage}[l]{0.22\columnwidth}
\vspace{6mm}
\includegraphics[width=1.1\columnwidth]{img/logo_soprahr_software.png}\\
\end{minipage}
\vskip1.5cm

\begin{center}
{\LARGE{\textbf{\textsc{Rapport de Projet de Fin d'\'Etudes}}}}\\
\vskip0.5cm
\large

{\textbf{Présenté en vue de l'obtention du}}\\
\vskip2mm
{\textbf{\@diplomaName}}\\
{\textbf{Spécialité : \@speciality}}\\
{}
\end{center}

\begin{center}
\textrm{Par}\\
\vskip0.3cm
{\ifthenelse{\boolean{isBinomal}}
    {% IF TRUE
        \begin{center}
            \large\textbf{\@author}~~~~~ et ~~~~~
            \large\textbf{\@secondAuthor}
        \end{center}
    }
    {\Large\textbf{\@author}}% FALSE
}
\vskip12mm

\definecolor{isiBlue}{RGB}{31, 78, 121}

\begin{changemargin}{-9mm}{0cm}
\begin{minipage}[l]{1.1\columnwidth}
\begin{tcolorbox}[colframe=isiBlue,colback=white,boxrule=0pt,toprule=3pt,bottomrule=3pt,arc=0pt,top=0mm,right=0mm,left=0mm,bottom=0mm,boxsep=0.5mm]{
    \begin{tcolorbox}[colframe=isiBlue,colback=white, boxrule=0pt,toprule=1pt,bottomrule=1pt,arc=0pt,enlarge bottom by=-0.9mm, auto outer arc]
        \centering
        {\huge\textbf{\@title}}
    \end{tcolorbox}
}
\end{tcolorbox}
\end{minipage}
\end{changemargin}

\end{center}
\vskip8mm%

\begin{center}
\large
% Trois minipages cote a cote ne s'alignaient pas : « Encadrant professionnel : »
% et « Sopra HR Software Tunisie » ne tenaient pas dans leur colonne et
% passaient a la ligne, ce qui rendait ces boites plus hautes que les autres.
% Comme elles etaient centrees verticalement ([c]), leurs lignes se retrouvaient
% plus haut que le libelle correspondant. Un tabular a trois colonnes de largeur
% naturelle supprime les deux problemes : chaque valeur est sur la meme ligne que
% son libelle, et rien ne revient a la ligne.
\begin{tabular}{@{}l@{\hspace{1.4em}}l@{\hspace{2.2em}}l@{}}
Encadrant professionnel : & \textbf{\@proFramerName} & \@proFramerSpeciality \\[5pt]
Encadrant académique : & \textbf{\@academicFramerName} & \@academicFramerSpeciality
\end{tabular}
\end{center}
\vskip16mm

\begin{center}
\large
Réalisé au sein de \@companyName\\
\vskip0.4cm
\begin{figure}[h]
\centering
{\color{isiBlue}{\fboxrule=2.5pt\fbox{\includegraphics[width=0.4\columnwidth]{logoSoprahr}}}}
\end{figure}
\end{center}

\begin{center}
\vskip0.4cm
    A.U : 2025-2026
\end{center}

\afterpage{\blankpage}